
\section{Dedução}
\label{sec:deducao}

A cinemática de um manipulador robótico permite estabelecer a relação entre as variáveis articulares e a posição e orientação do \textit{end effector}. No caso do MyCobot280, constituído por seis articulações rotativas, esta análise divide-se em dois problemas fundamentais: a cinemática direta e a cinemática inversa.

A cinemática direta permite determinar a posição e orientação do \textit{end effector} a partir dos ângulos das articulações. Por outro lado, a cinemática inversa consiste em determinar os ângulos articulares necessários para atingir uma determinada posição e orientação.

Nos capítulos seguintes, são apresentadas as formulações matemáticas utilizadas no desenvolvimento do projeto, que servem de base à implementação dos cálculos cinemáticos em MATLAB e à simulação dos movimentos do manipulador no CoppeliaSim.
